% Provenance (never printed): sources/book2/BRIEFS.md, g10-08. Willingness to pay, consumer and producer
% surplus, total surplus maximised at the competitive equilibrium, limits of the measure. Real numbers:
% figdata/grade-10/08-surplus-and-efficiency/ (Brynjolfsson, Collis and Eggers 2019, Table 1, cited rows).
% Brenland numbers invented.
\chapter{Surplus and Efficiency}\label{ch:g10:surplus-and-efficiency}

How much would someone have to pay you to give up search engines for a whole year: no web search on any device, for
twelve months? In 2017 a team of researchers asked this of thousands of Americans, offering amounts from one dollar to
twenty thousand. Half of them refused \usd{17530}. For email the median answer was \usd{8414}; for digital maps, \usd{3648}.
% ledger: g10:bce2019-search, g10:bce2019-email, g10:bce2019-maps, g10:bce2019-design

These services are free. Nobody pays for a web search, so the national accounts of the next chapters, which add up what is
paid for goods, count almost nothing for them. Yet the people who use them would give up a great deal to keep them. The gap
between what a good is worth to its buyer and what the buyer pays is called surplus. This chapter measures it, for buyers and
for sellers, and proves the central result about competitive markets: with the assumptions of the model, the equilibrium
quantity makes the \omterm{def:g10:surplus-and-efficiency:total-surplus}{total surplus} as large as it can be.

\begin{openingchart}{Median amount that American internet users would accept to give up a category of free digital goods for one
year, 2017, in dollars, on a times-ten scale.}{Brynjolfsson, Collis and Eggers (2019), Table 1}
\begin{tikzpicture}
\begin{axis}[omrecord, xbar, xmode=log, log basis x=10, log ticks with fixed point, log origin x=infty,
  width=0.78\linewidth, height=6.4cm, xmin=100, xmax=30000, bar width=7pt,
  ytick=data, yticklabels from table={figdata/grade-10/08-surplus-and-efficiency/out/digital_wta.csv}{category},
  y dir=reverse, ymin=0.4, ymax=8.6, ymajorgrids=false, xlabel={dollars a year (times-ten scale)},
  point meta=explicit, nodes near coords={\pgfmathprintnumber[fixed, precision=0, 1000 sep={\,}]{\pgfplotspointmeta}},
  nodes near coords style={font=\scriptsize, anchor=west}]
\addplot[fill=omDef!70, draw=omDef] table[x=wta2017, y=rank, meta=wta2017] {figdata/grade-10/08-surplus-and-efficiency/out/digital_wta.csv};
\end{axis}
\end{tikzpicture}
\end{openingchart}

\begin{recall}
From the grades 6--9 volume: a \emph{demand curve} gives the quantity buyers want at each price and a \emph{supply curve} the quantity
sellers offer. From \cref{ch:g10:supply}: the height of a supply curve at a quantity is the cost of the marginal unit. From
\cref{ch:g10:market-equilibrium}: the equilibrium of a competitive market, with price-taking buyers and sellers.
\end{recall}

\section{Willingness to pay and consumer surplus}\label{sec:g10:surplus-and-efficiency:consumer}

\begin{definition}[Willingness to pay]\label{def:g10:surplus-and-efficiency:wtp}
A buyer's \emph{willingness to pay}\index{willingness to pay} for a unit of a good is the largest amount of money the buyer would give for
it: at that price, buying and not buying are equally good to the buyer.
\end{definition}

Eight people in a Brenland town would each buy one ticket for a concert. Their \omterm{def:g10:surplus-and-efficiency:wtp}{willingness to pay}, in dalens, is 10, 9, 7, 6, 4, 3, 2 and 1.
Ranked from highest to lowest, these values form a staircase: at a price of 5 dalens, the first four buy; at 3.5, the first five do. The
staircase \emph{is} the demand curve of this small market: its height at the $k$-th ticket is the $k$-th highest \omterm{def:g10:surplus-and-efficiency:wtp}{willingness to pay}.

\begin{definition}[Consumer surplus]\label{def:g10:surplus-and-efficiency:consumer-surplus}
A buyer's \emph{consumer surplus}\index{consumer surplus} on a unit is its \omterm{def:g10:surplus-and-efficiency:wtp}{willingness to pay} minus the price paid. The consumer surplus of a
market is the sum over all units bought.
\end{definition}

At a price of 5 dalens, the four buyers gain $5 + 4 + 2 + 1 = 12$ dalens: the area of the staircase above the price line, up to the fourth
ticket.

\begin{omfigure}
\begin{tikzpicture}[econ, x=0.75cm, y=0.38cm]
\draw[econ axes] (0,11.5) node[left, font=\small] {dalens} |- (9,0) node[below left=12pt and 0pt, font=\small] {tickets};
\foreach \k/\v in {1/10,2/9,3/7,4/6} \fill[omDef!20] (\k-1,5) rectangle (\k,\v);
\draw[curve] (0,10) -- (1,10) -- (1,9) -- (2,9) -- (2,7) -- (3,7) -- (3,6) -- (4,6) -- (4,4) -- (5,4) -- (5,3) -- (6,3) -- (6,2) -- (7,2)
  -- (7,1) -- (8,1);
\draw[omThm, thick] (0,5) -- (8.5,5) node[right, font=\footnotesize] {price 5};
\foreach \k/\v in {1/10,2/9,3/7,4/6} \node[font=\scriptsize] at (\k-0.5,{(\v+5)/2}) {\the\numexpr\v-5\relax};
\foreach \y in {5,10} \draw (0,\y) -- (-0.15,\y) node[left, font=\footnotesize] {\y};
\foreach \x in {1,...,8} \node[below, font=\footnotesize] at (\x-0.5,0) {\x};
\end{tikzpicture}

{\small Eight buyers ranked by \omterm{def:g10:surplus-and-efficiency:wtp}{willingness to pay} form a staircase demand curve. At a price of 5 dalens the first four buy, and their \omterm{def:g10:surplus-and-efficiency:consumer-surplus}{consumer
surplus} is the shaded area, $5 + 4 + 2 + 1 = 12$ dalens.}
\end{omfigure}

With many buyers the steps become too small to see and the staircase becomes a smooth curve. \omterm{def:g10:surplus-and-efficiency:consumer-surplus}{Consumer surplus} is then the area between the
demand curve and the price line. For a straight demand line, it is a triangle.

\begin{example}[Consumer surplus under a line]\label{ex:g10:surplus-and-efficiency:cs-line}
For the demand $Q = 120 - 4p$, nobody buys above 30 dalens. At a price of 24, $Q = 24$ units are bought, and the \omterm{def:g10:surplus-and-efficiency:consumer-surplus}{consumer surplus} is the triangle
with base 24 and height $30 - 24 = 6$: $\tfrac12 \times 24 \times 6 = 72$ dalens.
\end{example}

\begin{measure}[Valuing a free good: asking people what they would accept]
A free good has no price, so no market reveals its value. The study of the opening chart asked people to choose between keeping a good and
receiving a sum of money for giving it up for a year. Each person saw one sum, drawn among nine amounts from \usd{1} to \usd{20000}; the share of
people who take the money at each amount traces a curve, and the median is the amount at which half accept. The study covered 19 digital goods and
6 others for comparison. Most of these choices were hypothetical. To check them, the researchers ran smaller experiments in which the choice was
binding: people who chose the money were paid and their use of the good was checked. In one, the median user needed about \usd{48} to give up
Facebook for one month. Two cautions remain: what people accept to give up a good is usually larger than what they would pay to get it, and an
online survey misses people who do not use the internet.
\end{measure}
% ledger: g10:bce2019-design, g10:bce2019-facebook

\section{Cost and producer surplus}\label{sec:g10:surplus-and-efficiency:producer}

On the other side, each seller has a lowest price at which it will sell a unit: the cost of producing it, including what the seller's time and
money could earn elsewhere.

\begin{definition}[Producer surplus]\label{def:g10:surplus-and-efficiency:producer-surplus}
A seller's \emph{producer surplus}\index{producer surplus} on a unit is the price received minus the cost of that unit. The producer surplus of a
market is the sum over all units sold: the area between the price line and the supply curve.
\end{definition}

For the supply $Q = 2p - 24$, nothing is offered below 12 dalens. At a price of 24, sellers offer 24 units, and the \omterm{def:g10:surplus-and-efficiency:producer-surplus}{producer surplus} is the triangle
with base 24 and height $24 - 12 = 12$: $\tfrac12 \times 24 \times 12 = 144$ dalens. \omterm{def:g10:surplus-and-efficiency:producer-surplus}{Producer surplus} is not profit: the cost of the fixed equipment
is not counted in it. It measures what sellers gain from selling, compared with not selling at all, once the equipment is there.

\section{Total surplus and the competitive equilibrium}\label{sec:g10:surplus-and-efficiency:total}

\begin{definition}[Total surplus]\label{def:g10:surplus-and-efficiency:total-surplus}
The \emph{total surplus}\index{total surplus} of a market is the sum of the consumer and \omterm{def:g10:surplus-and-efficiency:producer-surplus}{producer surplus}: for each unit traded, the buyer's
\omterm{def:g10:surplus-and-efficiency:wtp}{willingness to pay} minus the seller's cost. The price paid cancels out: it is taken from one side and given to the other.
\end{definition}

The demand $Q = 120 - 4p$ and the supply $Q = 2p - 24$ meet at $p^* = 24$ and $Q^* = 24$ (from $120 - 4p = 2p - 24$, $6p = 144$). \omterm{def:g10:surplus-and-efficiency:total-surplus}{Total surplus} is
$72 + 144 = 216$ dalens: the triangle between the two lines, from the axis to the equilibrium quantity, $\tfrac12 \times 24 \times (30 - 12) = 216$.

\begin{definition}[Allocative efficiency]\label{def:g10:surplus-and-efficiency:efficiency}
An allocation of a good is \emph{allocatively efficient}\index{allocative efficiency} when no change in who buys, who sells and how much is traded
could raise \omterm{def:g10:surplus-and-efficiency:total-surplus}{total surplus}.
\end{definition}

\begin{proposition}[The competitive equilibrium maximises total surplus]\label{prop:g10:surplus-and-efficiency:max}
In a market where each buyer wants at most one unit and each seller sells at most one, suppose the units are traded at the equilibrium price
of a competitive market, so that every buyer whose \omterm{def:g10:surplus-and-efficiency:wtp}{willingness to pay} is above the price buys and every seller whose cost is below the price
sells. Then no other choice of the units traded, of who buys them and of who sells them gives a larger \omterm{def:g10:surplus-and-efficiency:total-surplus}{total surplus}.
\end{proposition}
\begin{proof}
Rank the buyers by \omterm{def:g10:surplus-and-efficiency:wtp}{willingness to pay}, $v_1 \ge v_2 \ge \dots$, and the sellers by cost, $c_1 \le c_2 \le \dots$. If $q$ units are traded, \omterm{def:g10:surplus-and-efficiency:total-surplus}{total surplus}
is the sum of the buyers' values minus the sum of the sellers' costs. For a given $q$ it is largest when the buyers are the $q$ with the highest values and
the sellers the $q$ with the lowest costs: replacing a buyer by one who values the good more, or a seller by one with a lower cost, raises it. \omterm{def:g10:surplus-and-efficiency:total-surplus}{Total
surplus} with $q$ units is then $S(q) = (v_1 - c_1) + \dots + (v_q - c_q)$. Since the $v_k$ fall and the $c_k$ rise with $k$, the terms $v_k - c_k$ fall:
they are positive up to some $q^*$ and negative after. $S(q)$ is therefore largest at $q = q^*$, the last unit for which $v_k \ge c_k$. At the equilibrium
price $p^*$, the buyers with $v_k \ge p^*$ buy and the sellers with $c_k \le p^*$ sell: these are the $q^*$ buyers with the highest values and the $q^*$
sellers with the lowest costs, since $v_{q^*} \ge p^* \ge c_{q^*}$ and, for the next unit, $v_{q^*+1} < c_{q^*+1}$ leaves no price at which both would
trade.
\end{proof}

The proof makes no use of what anyone knows about the others: each buyer and seller only compares a value or a cost with the price. The price
does the sorting.

\begin{omfigure}
\begin{tikzpicture}[econ, x=0.13cm, y=0.16cm]
\begin{scope}
\draw[econ axes] (0,33) node[left, font=\small] {$p$} |- (52,0) node[below left=12pt and 0pt, font=\small] {$Q$};
\fill[omDef!20] (0,30) -- (24,24) -- (0,24) -- cycle;
\fill[omThm!15] (0,12) -- (24,24) -- (0,24) -- cycle;
\draw[demand] (0,30) -- (48,18);
\draw[supply] (0,12) -- (40,32);
\draw[guide] (24,24) -- (24,0);
\node[eqpoint] at (24,24) {};
\node[font=\scriptsize] at (7,26) {CS 72};
\node[font=\scriptsize] at (7,20.5) {PS 144};
\foreach \y in {12,24,30} \draw (0,\y) -- (-0.6,\y) node[left, font=\footnotesize] {\y};
\node[below, font=\footnotesize] at (24,0) {24};
\node[font=\scriptsize, anchor=north] at (24,-3.2) {at equilibrium};
\end{scope}
\begin{scope}[xshift=7.6cm]
\draw[econ axes] (0,33) node[left, font=\small] {$p$} |- (52,0) node[below left=12pt and 0pt, font=\small] {$Q$};
\fill[omDef!20] (0,30) -- (12,27) -- (12,18) -- (0,12) -- cycle;
\fill[gray!35] (12,27) -- (24,24) -- (12,18) -- cycle;
\draw[demand] (0,30) -- (48,18);
\draw[supply] (0,12) -- (40,32);
\draw[guide] (12,27) -- (12,0);
\node[font=\scriptsize] at (6,21) {162};
\node[font=\scriptsize, anchor=north west] at (26,17.5) {lost 54};
\draw[gray] (19,22.5) -- (26.3,17.3);
\foreach \y in {12,18,27} \draw (0,\y) -- (-0.6,\y) node[left, font=\footnotesize] {\y};
\node[below, font=\footnotesize] at (12,0) {12};
\node[font=\scriptsize, anchor=north] at (24,-3.2) {12 units only};
\end{scope}
\end{tikzpicture}

{\small Left: at the equilibrium (24 units at 24 dalens), \omterm{def:g10:surplus-and-efficiency:consumer-surplus}{consumer surplus} is 72 and \omterm{def:g10:surplus-and-efficiency:producer-surplus}{producer surplus} 144, a total of 216. Right: if only 12 units are
traded, the \omterm{def:g10:surplus-and-efficiency:total-surplus}{total surplus} is the area between the curves up to 12 units, 162; the grey triangle, 54, is lost.}
\end{omfigure}

The figure on the right shows what happens away from equilibrium. If only 12 units are traded, the twelfth buyer still values a unit at $(120 - 12)/4
= 27$ dalens and the twelfth seller's cost is $(12 + 24)/2 = 18$: there are units between the 12th and the 24th that buyers value more than they cost to
produce, and they are not traded. \omterm{def:g10:surplus-and-efficiency:total-surplus}{Total surplus} is the trapezoid $\tfrac12 \times (18 + 9) \times 12 = 162$, and $216 - 162 = 54$ dalens of surplus are
lost. Trading 36 units loses the same amount: units 25 to 36 cost more to produce than buyers value them. The lost area is called the deadweight loss;
\cref{ch:g10:what-prices-do} defines it.

\begin{model}{Efficiency of the competitive market}
\assumes Price-taking buyers and sellers who know the good's quality; nobody outside the market is affected by the trades; buyers' \omterm{def:g10:surplus-and-efficiency:wtp}{willingness to
  pay} and sellers' costs measure everything that matters to them about the good.
\predicts The equilibrium quantity is traded by the buyers who value the good most and the sellers who produce it most cheaply, and \omterm{def:g10:surplus-and-efficiency:total-surplus}{total surplus} is
  the largest possible.
\tested Markets run in laboratories, where the experimenter assigns values and costs and so knows the largest possible surplus, settle near the
  equilibrium price and quantity within a few rounds (Smith 1962), which by \cref{prop:g10:surplus-and-efficiency:max} is where \omterm{def:g10:surplus-and-efficiency:total-surplus}{total surplus} is largest.
\breaks When trades affect people outside the market: the sulphur dioxide from coal power stations harmed people who neither sold nor bought
  the power (grade 11). When sellers or buyers are few enough to move the price: a market with two or three sellers (grade 11). When buyers cannot see
  the quality of what they buy: used cars (grade 11).
\end{model}
% ledger: g10:smith-1962

\section{What the result does not say}\label{sec:g10:surplus-and-efficiency:limits}

The proposition is precise, and it is easy to read more into it than it says.

\begin{itemize}
\item \textbf{It says nothing about who gets the surplus.} In the example above sellers get 144 and buyers 72. A different split of the same 216 would be
  just as efficient. Efficiency and the division of the gains are separate questions.
\item \textbf{It counts only the people in the market.} If a factory's smoke harms the neighbours, their loss is in no buyer's value and no seller's cost.
  The market can then trade too much, and the proposition does not apply.
\item \textbf{It measures in money.} \omterm{def:g10:surplus-and-efficiency:wtp}{Willingness to pay} depends on what one has. A rich person will pay more for a concert ticket than a poor person who
  would enjoy it more. \omterm{def:g10:surplus-and-efficiency:total-surplus}{Total surplus} adds dalens, not happiness: the equilibrium gives the ticket to whoever will pay most, and that is the person who
  values it most only in money.
\item \textbf{It is about one market at a time.} A change in one market moves others, as \cref{ch:g10:market-equilibrium} showed, and the surplus
  of one market does not count those effects.
\end{itemize}

The opening chart shows the measure's reach and its limit together. The value that users place on free services is surplus that no account of
production records; but it is measured by what people say they would accept, and two people with the same answer need not gain the same thing from the
service.

\section{Exercises}

\begin{exercise}[$\star$]\label{exo:g10:surplus-and-efficiency:1}
Ana is willing to pay 40 dalens for a pair of shoes and buys them for 32. What is her \omterm{def:g10:surplus-and-efficiency:consumer-surplus}{consumer surplus}? What would it be if the price were 45?
\end{exercise}

\begin{exercise}[$\star$]\label{exo:g10:surplus-and-efficiency:2}
A baker's cost for a cake is 12 dalens and he sells it for 20. What is his \omterm{def:g10:surplus-and-efficiency:producer-surplus}{producer surplus}? Below what price would he not sell?
\end{exercise}

\begin{exercise}[$\star$]\label{exo:g10:surplus-and-efficiency:3}
Read the opening chart. Which category had the largest median? About how many times larger was it than the median for social media? Which two
categories were closest?
\end{exercise}

\begin{exercise}[$\star$]\label{exo:g10:surplus-and-efficiency:4}
The demand for a good is $Q = 60 - 2p$. Compute the \omterm{def:g10:surplus-and-efficiency:consumer-surplus}{consumer surplus} at a price of 10 dalens.
\end{exercise}

\begin{exercise}[$\star$]\label{exo:g10:surplus-and-efficiency:5}
The supply of a good is $Q = 3p - 15$, for $p \ge 5$. Compute the \omterm{def:g10:surplus-and-efficiency:producer-surplus}{producer surplus} at a price of 15 dalens.
\end{exercise}

\begin{exercise}[$\star\star$]\label{exo:g10:surplus-and-efficiency:6}
Demand is $Q = 100 - 2p$ and supply $Q = 3p - 60$, for $p \ge 20$. Find the equilibrium and the consumer, producer and \omterm{def:g10:surplus-and-efficiency:total-surplus}{total surplus}.
\end{exercise}

\begin{exercise}[$\star\star$]\label{exo:g10:surplus-and-efficiency:7}
Demand is $Q = 80 - p$ and supply $Q = p$. Find the equilibrium and the \omterm{def:g10:surplus-and-efficiency:total-surplus}{total surplus}. Show that buyers and sellers get the same surplus,
and explain why.
\end{exercise}

\begin{exercise}[$\star\star$]\label{exo:g10:surplus-and-efficiency:8}
In the market of \cref{exo:g10:surplus-and-efficiency:6}, compute the \omterm{def:g10:surplus-and-efficiency:total-surplus}{total surplus} if only 18 units are traded (by the buyers who value them most and the
sellers with the lowest costs), and the surplus lost.
\end{exercise}

\begin{exercise}[$\star\star$]\label{exo:g10:surplus-and-efficiency:9}
In the market $Q^d = 120 - 4p$, $Q^s = 2p - 24$ of the text, a new technique lowers every seller's cost by 3 dalens, and supply becomes $Q = 2p - 18$.
Find the new equilibrium and the new consumer and \omterm{def:g10:surplus-and-efficiency:producer-surplus}{producer surplus}. Who gains?
\end{exercise}

\begin{exercise}[$\star\star$]\label{exo:g10:surplus-and-efficiency:10}
In the concert example, a ninth buyer arrives with a \omterm{def:g10:surplus-and-efficiency:wtp}{willingness to pay} of 8 dalens. Draw the new staircase. At a price of 5, who buys, and what is the
\omterm{def:g10:surplus-and-efficiency:consumer-surplus}{consumer surplus}?
\end{exercise}

\begin{exercise}[$\star\star$]\label{exo:g10:surplus-and-efficiency:11}
In the market of the text, demand shifts to $Q = 126 - 4p$. Find the new equilibrium and compute the consumer and \omterm{def:g10:surplus-and-efficiency:producer-surplus}{producer surplus}. Which side gained
more?
\end{exercise}

\begin{exercise}[$\star\star\star$]\label{exo:g10:surplus-and-efficiency:12}
Describe evidence that would show that a real market fails to trade units that buyers value more than they cost to produce. What would you need to know,
and how could you learn it?
\end{exercise}

\begin{exercise}[$\star\star\star$]\label{exo:g10:surplus-and-efficiency:13}
For the demand $Q = a - bp$ and the supply $Q = c + dp$ (with $c < 0 < a$, $b, d > 0$), let $Q < Q^*$ be a quantity traded by the buyers with the highest values
and the sellers with the lowest costs. Using only the areas of triangles, show that the lost surplus is $\tfrac12 (Q^* - Q)\,g$, where $g$ is the gap
between the demand price and the supply price at $Q$, and deduce that \omterm{def:g10:surplus-and-efficiency:total-surplus}{total surplus} is largest at $Q^*$.
\end{exercise}

\begin{exercise}[$\star\star\star$]\label{exo:g10:surplus-and-efficiency:14}
Give an example in which moving a unit of a good from one person to another lowers \omterm{def:g10:surplus-and-efficiency:total-surplus}{total surplus} but, plausibly, makes the two people together better off.
What does this say about using \omterm{def:g10:surplus-and-efficiency:total-surplus}{total surplus} to judge a market?
\end{exercise}

\begin{exercise}[$\star\star\star$]\label{exo:g10:surplus-and-efficiency:15}
A farm sells eggs in a competitive market, and its smell costs each of its 50 neighbours 2 dalens per hundred eggs produced. Explain why the market
equilibrium no longer maximises the surplus of everyone concerned. Would the right quantity of eggs be larger or smaller?
\end{exercise}

\section{Problem: The Saturday Market}

\begin{problem}[{Weekend problem --- ten buyers, eight sellers and a cap on stalls}]
\label{pb:g10:surplus-and-efficiency:1}
Every Saturday, in a square of a Brenland town, eight farmers can each sell one basket of vegetables and ten shoppers can each buy one. The shoppers'
\omterm{def:g10:surplus-and-efficiency:wtp}{willingness to pay} and the farmers' costs, in dalens, are:
\begin{center}
\small
\begin{tabular}{l c c c c c c c c c c}
\toprule
shoppers & 30 & 28 & 25 & 22 & 20 & 18 & 15 & 12 & 10 & 8 \\
farmers  & 6 & 8 & 11 & 13 & 16 & 19 & 23 & 26 & & \\
\bottomrule
\end{tabular}
\end{center}

\medskip
\noindent\textbf{Part I --- The market without rules.}
\begin{enumerate}
\item Draw the staircase demand and staircase supply on one graph.
\item How many baskets are wanted at a price of 21 dalens, and how many offered? At 17?
\item Find the range of prices at which the number of baskets wanted equals the number offered. How many baskets are traded?
\item Which shoppers buy and which farmers sell?
\item Compute the largest possible \omterm{def:g10:surplus-and-efficiency:total-surplus}{total surplus}, and check that the market reaches it.
\end{enumerate}

\medskip
\noindent\textbf{Part II --- Splitting the surplus.}
\begin{enumerate}[resume]
\item At a price of 18.5 dalens, compute the \omterm{def:g10:surplus-and-efficiency:consumer-surplus}{consumer surplus} and the \omterm{def:g10:surplus-and-efficiency:producer-surplus}{producer surplus}.
\item At a price of 18 and at a price of 19, compute the two surpluses again. Does the total change?
\item The sixth shopper (18) and the sixth farmer (19) do not trade. Why would a trade between them lower \omterm{def:g10:surplus-and-efficiency:total-surplus}{total surplus}?
\item The shopper who values a basket at 30 pays the same price as the one who values it at 20. Who gains more surplus? Is the market less efficient
  because of this?
\end{enumerate}

\medskip
\noindent\textbf{Part III --- A cap on stalls.}
\begin{enumerate}[resume]
\item The town limits the square to 3 stalls, given to the three farmers with the lowest costs. How many baskets are traded now?
\item Between which prices does the market now settle? (Find the prices at which exactly three shoppers want a basket.)
\item Compute the \omterm{def:g10:surplus-and-efficiency:total-surplus}{total surplus} with the cap.
\item How much surplus is lost, in dalens and as a share of the largest possible surplus?
\item Which farmers and which shoppers lose? Can any group gain from the cap?
\end{enumerate}

\medskip
\noindent\textbf{Part IV --- Who gets the stalls.}
\begin{enumerate}[resume]
\item Suppose instead that the 3 stalls go by lottery, to the farmers with costs 11, 19 and 26. Between which prices does the market settle now, how many
  baskets are traded, and what is the \omterm{def:g10:surplus-and-efficiency:total-surplus}{total surplus}?
\item How much is lost now, in dalens and as a share of the largest possible surplus?
\item Why is the loss larger than in question 13?
\item With the lottery, the farmer with a cost of 6 has no stall but the farmer with a cost of 26 has one, which he does not use. Explain why both would agree
  to hand the stall over for a payment.
\item After that trade the stalls are held by the farmers with costs 6, 11 and 19. How many baskets are traded, what is the \omterm{def:g10:surplus-and-efficiency:total-surplus}{total surplus}, and how much of
  the lottery's loss has the trade recovered?
\item State the result of the problem: the surplus lost to the stall cap, in dalens and as a share of the maximum.
\end{enumerate}
\end{problem}
